\chapter[Terminal, naturalidad y l\'imite]{Terminal multivaluado, naturalidad \'unica y paso al l\'imite}
\label{chap:pdfv2-terminal-naturalidad-limite}

\begin{thesisbox}
El continuo discreto produce un solo sistema de historias, extensiones,
memorias y terminales. Los discriminantes internos
\(\mathrm{sp},\mathrm{sol},\mathrm{gau},\mathrm{coh},\mathrm{det}\)
nombran cinco caracter\'isticas simult\'aneas de ese mismo sistema: no fijan
entradas distintas ni autorizan a escoger historias diferentes. En este
cap\'itulo el terminal conserva todas las continuaciones, la conexi\'on
non\'adica permanece como correspondencia y la naturalidad se expresa por una
sola identidad.
\end{thesisbox}

\section{Historias supervivientes con testigos de extensi\'on}

Sea
\(\widehat x_k\in\widehat{\mathcal X}^{\rm tot}_k\) el estado total de
\eqref{eq:continuo-comun-dominio-total}. Para \(N\geq k\), definimos sin
otro predicado de pertenencia
\begin{equation}
\label{eq:pdfv2-terminal-historias-finita}
 \mathscr H_{k,N}(\widehat x_k)
 :=\operatorname{Term}_{k,N}(\widehat x_k).
\end{equation}
Cada
\(h=(\varepsilon_{k+1},\ldots,\varepsilon_N)\) de la derecha contiene por
\eqref{eq:continuo-comun-vertice-arbol} sus estados intermedios, los testigos
de \(\operatorname{Ext}\), los carries, las fronteras, las rutas y el ledger.
Cada arista satisface la acción de fase
\eqref{eq:continuo-comun-accion-fase-tpk}; por ello todo subcamino de nueve
pasos determina, con sus nueve emisiones y no sólo con su longitud, un
elemento de \(\mathsf Z^\Gamma_{9,j}\) y un testigo de
\(\Gamma_{9,j}\) por
\eqref{eq:continuo-comun-gamma9-apice}--
\eqref{eq:continuo-comun-gamma9-relacion}. Por tanto, Ext y la holonomía
non\'adica nacen de la misma palabra fase--compatible; no se añade a
\eqref{eq:pdfv2-terminal-historias-finita} un filtro de supervivencia ni una
segunda condici\'on \(\Gamma_9\).

Para \(k\leq\ell\leq N\), sea
\begin{equation}
\label{eq:pdfv2-terminal-prefijos-intermedios}
 \mathscr P_{\ell\mid k}(\widehat x_k)
 :=\{h_{[k,\ell]}:h\in\mathscr H_{k,N}(\widehat x_k)
       \text{ para alg\'un }N\geq\ell\}.
\end{equation}
Para \(p\in\mathscr P_{\ell\mid k}(\widehat x_k)\), usamos la abreviatura
\begin{equation}
\label{eq:pdfv2-terminal-continuaciones-prefijo}
 \mathscr H_{\ell,N}(p)
 :=\{q:p\star q\in\mathscr H_{k,N}(\widehat x_k)\}.
\end{equation}
La prolongaci\'on neutra de \eqref{eq:continuo-comun-seccion-total} extiende
cada prefijo a cualquier horizonte. En consecuencia,
\begin{equation}
\label{eq:pdfv2-terminal-no-vacuidad-historias}
 \mathscr H_{k,N}(\widehat x_k)\neq\varnothing,
 \qquad
 \mathscr P_{\ell\mid k}(\widehat x_k)\neq\varnothing,
\end{equation}
por \cref{prop:continuo-comun-arbol-finito-no-vacio}.

El truncamiento
\begin{equation}
\label{eq:pdfv2-terminal-truncamiento-historias}
 \pi_{N+1,N}^{\mathscr H}:
 \mathscr H_{k,N+1}(\widehat x_k)
 \longrightarrow\mathscr H_{k,N}(\widehat x_k)
\end{equation}
es exactamente \eqref{eq:continuo-comun-proyeccion-terminal} y es
sobreyectivo por \cref{lem:continuo-comun-terminal-sobreyectivo}. La
supervivencia es, por
\eqref{eq:continuo-comun-surv-igual-total}, una salida demostrada del dominio
total y no una selecci\'on previa.

\section{Registro terminal generado y sus trece coordenadas}

Fijemos una historia
\[
 h=(\varepsilon_{j+1}:x_j\longrightarrow x_{j+1})_{j=k}^{N-1},
 \qquad
 \gamma_{k,k}:=I_{x_k},\qquad
 \gamma_{k,j+1}:=\gamma_{k,j}\star\varepsilon_{j+1},
\]
La notación operatorial
\(U_{\gamma_{k,j+1}}
=U_{\varepsilon_{j+1}}U_{\gamma_{k,j}}\) conserva el orden de aplicación,
mientras \(\star\) conserva el orden cronológico de la palabra. Fijamos
además un nivel de coeficientes \(r\geq1\). Para abreviar, escribimos
\[
 C_{x_j,r}:=C_{j,r}(x_j),\qquad
 F_{\varepsilon_{j+1}}:C_{x_{j+1},r}\longrightarrow C_{x_j,r}
\]
para el complejo y el mapa de complejos de
\eqref{eq:pdfv2-cor-cono-celda} y
\eqref{eq:pdfv2-cor-refinamiento-celda}. El tipo del registro de nivel
\(j\), denotado \(\mathfrak R_{j,r}\), es un registro dependiente con los
trece campos siguientes. En cada campo indicamos el tipo, la semilla y la
acción de una arista \(\varepsilon:=\varepsilon_{j+1}\).

El estado \(x_k\) ya contiene la palabra completa
\(\gamma_{0,k}=\varepsilon_1\star\cdots\star\varepsilon_k\). Por ello una
``semilla en \(k\)'' nunca reinicia la genealogía: es el fold de las mismas
actualizaciones desde la raíz hasta \(k\). Denotamos
\begin{equation}
\label{eq:pdfv2-terminal-prefijo-heredado}
 \begin{gathered}
 M_k^\pm:=\sum_{i=1}^{k}b_{\varepsilon_i}^\pm,qquad
 D_k:=\sum_{i=1}^{k}b_{\varepsilon_i},qquad
 H_k:=\sum_{i=1}^{k}|b_{\varepsilon_i}|,\\
 \mathbf c_{\leq k}:=
  \bigl(c_r(\varepsilon_i,\gamma_{0,i-1})\bigr)_{1\leq i\leq k},
 \qquad
 \mathbf k_{\leq k}:=
  \bigl(\kappa_{i-1}^{\rm can}\bigr)_{1\leq i\leq k},\\
 \mathbf F_{\leq k}:=
  \bigl(F_{\varepsilon_i}:C_{x_i,r}\to C_{x_{i-1},r}\bigr)_{1\leq i\leq k}.
 \end{gathered}
\end{equation}
Todas las listas vacías que aparecen abajo se entienden literalmente sólo
en \(k=0\); para \(k>0\) se sustituyen por los prefijos heredados de
\eqref{eq:pdfv2-terminal-prefijo-heredado} y del ledger de
\(\gamma_{0,k}\). Así el registro no conserva una copia opaca de \(x_k\),
pero tampoco pierde su pasado APP--TRIT--TPK.

\paragraph{1. Fibra \(\mathsf F\).}
La fibra elemental es
\begin{equation}
\label{eq:pdfv2-terminal-tipo-fibra}
 \operatorname{Fib}_r(x):=
 \bigl(\mathsf H_x^\Sigma\oplus\mathsf H_x^\Pi,\,
       \mathsf M_x,\,
       \operatorname{Cell}_{\rm TPK}(x),\,C_{x,r}\bigr).
\end{equation}
El campo \(\mathsf F_{k,j}\) es una palabra en la categoría de estas fibras.
Su semilla y actualización son
\begin{align}
 \mathsf F_{k,k}
 &:=\bigl(\operatorname{Fib}_r(x_0),I\bigr)\star
 \mathop{\star}_{i=1}^{k}
 \bigl(\operatorname{Fib}_r(x_i),U_{\varepsilon_i},
 \mathsf C_{\varepsilon_i},\operatorname{Cell}_{\rm TPK}(\varepsilon_i),
 F_{\varepsilon_i}\bigr),
 \label{eq:pdfv2-terminal-semilla-f}\\
 \operatorname{Upd}^{\mathsf F}_{\varepsilon}\mathsf F_{k,j}
 &:=
 \mathsf F_{k,j}\star
 \bigl(\operatorname{Fib}_r(x_{j+1}),U_\varepsilon,
 \mathsf C_\varepsilon,\operatorname{Cell}_{\rm TPK}(\varepsilon),
 F_\varepsilon\bigr).
 \label{eq:pdfv2-terminal-update-f}
\end{align}
Es no vacía porque las dos hojas, la memoria, la unidad celular
\eqref{eq:pdfv2-cor-cell-unidad} y el complejo de puertos ya están
construidos. La composición de sus cinco mapas prueba la naturalidad.

\paragraph{2. Extensiones \(\mathsf E\).}
El codominio es la categoría de palabras de relaciones tipadas. Ponemos
\begin{align}
 \mathsf E_{k,k}
 &:=(I_{x_0},\operatorname{Ext}_{I_{x_0}})\star
 \mathop{\star}_{i=1}^{k}
 \bigl(\varepsilon_i,\operatorname{Ext}_{\varepsilon_i},
 \operatorname{Ext}_{\varepsilon_i}\circ
 \operatorname{Ext}_{\gamma_{0,i-1}}
 \subseteq\operatorname{Ext}_{\gamma_{0,i}}\bigr),
 \label{eq:pdfv2-terminal-semilla-e}\\
 \operatorname{Upd}^{\mathsf E}_{\varepsilon}\mathsf E_{k,j}
 &:=
 \mathsf E_{k,j}\star
 \bigl(\varepsilon,\operatorname{Ext}_{\varepsilon},
 \operatorname{Ext}_{\varepsilon}\circ
 \operatorname{Ext}_{\gamma_{0,j}}
 \subseteq\operatorname{Ext}_{\gamma_{0,j+1}}\bigr).
 \label{eq:pdfv2-terminal-update-e}
\end{align}
Identidad y emisión neutra prueban no vacuidad; asociatividad de la
composición relacional prueba naturalidad.

\paragraph{3. Relaciones \(\mathsf R\).}
Sea \(\operatorname{RelBlock}_{j,r}\) el tipo de bloques de ecuaciones
etiquetados por la arista y el prefijo sobre los que se verifican, y sea
\(\operatorname{Pres}_{j,r}:=\operatorname{List}
(\operatorname{RelBlock}_{j,r})\). El bloque unidad y la semilla son
\begin{equation}
\label{eq:pdfv2-terminal-semilla-r}
 \mathsf B_I:=(I_{x_0};U_I=I,\ \mathsf U_I=I,\ \kappa(I)=0,\ \partial I=0),
 \qquad
 \mathsf R_{k,k}:=(\mathsf B_I)\star
 \mathop{\star}_{i=1}^{k}
 \mathsf B_R(\varepsilon_i;\gamma_{0,i-1}).
\end{equation}
Para una arista \(\varepsilon:x_j\to x_{j+1}\), el bloque nuevo es la
tupla etiquetada
\begin{equation}
\label{eq:pdfv2-terminal-bloque-r}
 \begin{split}
 \mathsf B_R(\varepsilon;\gamma_{0,j}):=\bigl(
 &\varepsilon,\gamma_{0,j};\ 
 \delta^\diamond(\varepsilon)
 =\operatorname{res}^\diamond(\varepsilon)
  +9\operatorname{quo}^\diamond(\varepsilon),\\
 &\Delta(\varepsilon)=\operatorname{or}(\varepsilon)
  +3\kappa(\varepsilon),\quad
 U_{\gamma_{0,j+1}}=U_\varepsilon U_{\gamma_{0,j}},\\
 &\mathsf U_{\gamma_{0,j+1}}
 =\mathsf U_\varepsilon\mathsf U_{\gamma_{0,j}},\quad
 \kappa(\gamma_{0,j+1})
 =\kappa(\varepsilon)+\mathsf C_\varepsilon\kappa(\gamma_{0,j}),\\
 &\partial\gamma_{0,j+1}
 =\partial\varepsilon+\partial\gamma_{0,j},\quad
 \operatorname{Cpl}_\varepsilon
 \text{ de }\eqref{eq:pdfv2-cor-acoplamientos-unicos}
 \bigr).
 \end{split}
\end{equation}
La actualización concatena el bloque sin perder fecha ni multiplicidad:
\begin{equation}
\label{eq:pdfv2-terminal-update-r}
 \operatorname{Upd}^{\mathsf R}_{\varepsilon}\mathsf R_{k,j}
 :=\mathsf R_{k,j}\star\mathsf B_R(\varepsilon;\gamma_{0,j}).
\end{equation}
El truncamiento \(\operatorname{tr}^{\mathsf R}_{j+1,j}\) elimina el último
bloque etiquetado. Además
\(a_*\mathsf B_R(\varepsilon;\gamma)
=\mathsf B_R(a\varepsilon;a\gamma)\), porque cada ecuación es natural en
la misma arista y el mismo prefijo.

\paragraph{4. Números \(\mathsf N\).}
Para una arista \(\varepsilon:x_j\to x_{j+1}\), definimos el dato numérico
\begin{equation}
\label{eq:pdfv2-terminal-numero-arista}
 \nu(\varepsilon):=
 \bigl(j+1,\mathsf Q_{j+1},g_{j+1},q^{(9)}_{j+1},\beta_{j+1},
 \operatorname{res}^{\Sigma},\operatorname{quo}^{\Sigma},
 \operatorname{res}^{\Pi},\operatorname{quo}^{\Pi},
 \operatorname{or},\kappa,\Delta\bigr)(\varepsilon).
\end{equation}
Su tipo es
\[
 \mathbb N\times\widetilde{\mathcal O}_9\times C_9\times\mathbb N
 \times\mathbb Z/12\mathbb Z
 \times\{0,\ldots,8\}^2\times\mathbb Z^5.
\]
La semilla contiene la lista numérica completa del prefijo \(0\to k\);
la actualización añade exactamente el dato de la arista siguiente:
\begin{equation}
\label{eq:pdfv2-terminal-update-n}
 \mathsf N_{k,k}:=
 ((0,\mathsf Q_0,g_0,q^{(9)}_0,\beta_0;0,\ldots,0))
 \star\mathop{\star}_{i=1}^{k}\nu(\varepsilon_i),\qquad
 \operatorname{Upd}^{\mathsf N}_{\varepsilon}\mathsf N_{k,j}
 :=\mathsf N_{k,j}\star\nu(\varepsilon).
\end{equation}
Totalidad procede de las divisiones APP, la división TRIT y el odómetro TPK.

\paragraph{5. Orientación \(\mathsf O\).}
El tipo de un estado de orientación es
\[
 \mathbb Z\times
 (\mathsf H^\Sigma_x\oplus\mathsf H^\Pi_x)
 \times\mathbb Z[\mathsf A],
\]
con semilla y actualización
\begin{align}
 \mathsf O_{k,k}&:=
 \left(\sum_{i=1}^{k}\operatorname{or}(\varepsilon_i),
       \operatorname{Sh}(x_k),\partial\gamma_{0,k}\right),
 \label{eq:pdfv2-terminal-semilla-o}\\
 \operatorname{Upd}^{\mathsf O}_{\varepsilon}
 (\mathfrak o,\operatorname{Sh},\mathfrak b)
 &:=
 \bigl(\mathfrak o+\operatorname{or}(\varepsilon),
 U_\varepsilon\operatorname{Sh},
 \mathfrak b+\partial\varepsilon\bigr).
 \label{eq:pdfv2-terminal-update-o}
\end{align}
Las hojas APP prueban no vacuidad. Functorialidad de \(U\), aditividad de
\(\operatorname{or}\) y telescopía de la frontera prueban naturalidad.

\paragraph{6. Carry \(\mathsf K\).}
Usamos los lifts enteros tipados en
\eqref{eq:pdfv2-cor-carry-matricial-def}. El campo tiene tipo
\[
 \mathsf M_{x_j}\times
 \operatorname{Hom}(\mathsf M_{x_0},\mathsf M_{x_j})
 \times\operatorname{Mat}_5(\mathbb Z)
 \times\operatorname{List}(\operatorname{Mat}_5(\mathbb Z))
 \times\operatorname{List}(\mathbb Q_{\geq0}).
\]
Para \(\gamma_{k,j}=\varepsilon_j\cdots\varepsilon_{k+1}\), fijamos
explícitamente
\begin{equation}
\label{eq:pdfv2-terminal-lift-entero-ruta}
 \widetilde L_{\gamma_{k,k}}:=I_5,\qquad
 \widetilde L_{\gamma_{k,j}}
 :=\widetilde L_r(\varepsilon_j)\cdots
   \widetilde L_r(\varepsilon_{k+1}).
\end{equation}
Para el prefijo heredado fijamos asimismo
\[
 \widetilde L_{\gamma_{0,k}}
 :=\widetilde L_r(\varepsilon_k)\cdots
   \widetilde L_r(\varepsilon_1).
\]
La semilla y actualización son
\begin{align}
 \mathsf K_{k,k}&:=
 \bigl(\kappa(\gamma_{0,k}),\mathsf C_{\gamma_{0,k}},
       \widetilde L_{\gamma_{0,k}},
       \mathbf c_{\leq k},\mathbf k_{\leq k}\bigr),
 \label{eq:pdfv2-terminal-semilla-k}\\
 \operatorname{Upd}^{\mathsf K}_{\varepsilon}
 (\kappa_\gamma,\mathsf C_\gamma,\widetilde L_\gamma,
   \mathbf c,\mathbf k)
 &:=
 \bigl(\kappa(\varepsilon)+\mathsf C_\varepsilon\kappa_\gamma,\,
 \mathsf C_\varepsilon\mathsf C_\gamma,\,
 \widetilde L_r(\varepsilon)\widetilde L_\gamma,\,
 \mathbf c\star c_r(\varepsilon,\gamma),\,
 \mathbf k\star\kappa_j^{\rm can}\bigr).
 \label{eq:pdfv2-terminal-update-k}
\end{align}
La identidad proporciona la semilla; los dos cociclos
\eqref{eq:continuo-comun-cociclo-carry} y
\eqref{eq:pdfv2-cor-cociclo-matricial} prueban independencia del
parentizado y naturalidad.

\paragraph{7. Memoria \(\mathsf M\).}
Su tipo conserva a la vez memoria afín, ledger y memorias firmadas:
\begin{equation}
\label{eq:pdfv2-terminal-tipo-ledger-entry}
 \operatorname{LedgerEntry}:=
 \{\lambda(\varepsilon):\varepsilon\text{ emisión elemental tipada}\}.
\end{equation}
\[
 \mathsf M_x\times\operatorname{List}(\operatorname{LedgerEntry})
 \times\mathbb N^2\times\mathbb Z\times\mathbb N.
\]
Ponemos
\begin{align}
 \mathsf M_{k,k}&:=
 \bigl(m_k,\operatorname{Led}(\gamma_{0,k}),
       M_k^+,M_k^-,D_k,H_k\bigr),
 \label{eq:pdfv2-terminal-semilla-m}\\
 \operatorname{Upd}^{\mathsf M}_{\varepsilon}
 (m,\mathbf{Led},M^+,M^-,D,H)
 &:=
 \bigl(\mathsf C_\varepsilon m+\kappa(\varepsilon),\,
 \mathbf{Led}\star\lambda(\varepsilon),\,
 M^++b_\varepsilon^+,\,M^-+b_\varepsilon^-,\,
 D+b_\varepsilon,\,H+|b_\varepsilon|\bigr).
 \label{eq:pdfv2-terminal-update-m}
\end{align}
La acción afín y \eqref{eq:pdfv2-cor-memorias} prueban totalidad y
naturalidad; el pasado no se recalcula del último valor.

\paragraph{8. Frontera \(\mathsf B\).}
El tipo es una cadena de frontera acompañada por una palabra de complejos
basados y mapas de cadena. Su semilla y actualización son
\[
 \mathsf B_{k,j}\in
 \operatorname{Rec}[\mathfrak b\in\mathbb Z[\mathsf A_j];
 C_{x_j,r}\in\operatorname{Ch}_{\rm bas};
 \mathbf F\in\operatorname{Path}_{\operatorname{Ch}_{\rm bas}}
 (C_{x_j,r},C_{x_0,r})].
\]
\begin{align}
 \mathsf B_{k,k}&:=
 \bigl(\partial\gamma_{0,k},C_{x_k,r},\mathbf F_{\leq k}\bigr),
 \label{eq:pdfv2-terminal-semilla-b}\\
 \operatorname{Upd}^{\mathsf B}_{\varepsilon}
 (\mathfrak b,C_{x_j,r},\mathbf F)
 &:=
 \bigl(\mathfrak b+\partial\varepsilon,\,
 C_{x_{j+1},r},\,\mathbf F\star F_\varepsilon\bigr).
 \label{eq:pdfv2-terminal-update-b}
\end{align}
No vacuidad procede de los puertos comunes; \(\partial^2=0\),
\(F_{\beta\varepsilon}=F_\varepsilon F_\beta\) y
\eqref{eq:continuo-comun-frontera-telescopica} prueban compatibilidad.

\paragraph{9. Ruta \(\mathsf G\).}
El campo es la lista de prefijos en la categoría de caminos:
\begin{equation}
\label{eq:pdfv2-terminal-update-g}
 \mathsf G_{k,k}:=(\gamma_{0,i})_{0\leq i\leq k},\qquad
 \operatorname{Upd}^{\mathsf G}_{\varepsilon}\mathsf G_{k,j}
 :=\mathsf G_{k,j}\star\gamma_{0,j+1}.
\end{equation}
Identidad, concatenación asociativa y
\(a(\varepsilon\gamma)=a(\varepsilon)a(\gamma)\) prueban las tres
propiedades requeridas.

\paragraph{10. Multisección \(\mathsf X\).}
Para \(M\geq j\), definimos por acción
\begin{equation}
\label{eq:pdfv2-terminal-multiseccion-interna}
 \mathsf X_{k,j}(M):=
 \{\gamma_{k,j}\star q:
 q\in\operatorname{Term}_{j,M}(x_j)\}.
\end{equation}
Así
\[
 \mathsf X_{k,k}(M)=\operatorname{Term}_{k,M}(x_k),
\]
y
\begin{equation}
\label{eq:pdfv2-terminal-update-x}
 \operatorname{Upd}^{\mathsf X}_{\varepsilon}\mathsf X_{k,j}(M)
 :=
 \{\gamma_{k,j}\star\varepsilon\star q:
 q\in\operatorname{Term}_{j+1,M}(x_{j+1})\}.
\end{equation}
La prolongación neutra prueba no vacuidad. Borrar el último paso y actuar
por una simetría conmutan literalmente con la concatenación.

\paragraph{11. Supervivencia \(\mathsf S\).}
Aquí no se guarda una proposición, sino testigos. Sea
\begin{equation}
\label{eq:pdfv2-terminal-testigo-neutro}
 s_{j,j}:=I_{x_j},\qquad
 s_{j,M+1}:=s_{j,M}\star
 \varepsilon_M^0(t(s_{j,M})).
\end{equation}
La semilla es
\(\mathsf S_{k,k}:=(s_{k,M})_{M\geq k}\), y
\begin{equation}
\label{eq:pdfv2-terminal-update-s}
 \operatorname{Upd}^{\mathsf S}_{\varepsilon}\mathsf S_{k,j}
 :=(\gamma_{k,j}\star\varepsilon\star s_{j+1,M})_{M\geq j+1}.
\end{equation}
Por construcción
\(\pi_{M+1,M}s_{j,M+1}=s_{j,M}\); la naturalidad de la emisión neutra
prueba la equivariancia.

\paragraph{12. Terminal \(\mathsf T\).}
El terminal es el sistema inverso, no una cola escogida:
\begin{equation}
\label{eq:pdfv2-terminal-sistema-inverso-campo}
 \mathsf T(x_j):=
 \bigl((\operatorname{Term}_{j,M}(x_j))_{M\geq j},
       (\pi_{M+1,M})_{M\geq j}\bigr).
\end{equation}
Para cada \(M\geq j\), el mapa de prefijo tipado es
\begin{equation}
\label{eq:pdfv2-terminal-update-t}
 \operatorname{pref}_{\gamma_{k,j},M}:
 \operatorname{Term}_{j,M}(x_j)\longrightarrow
 \operatorname{Term}_{k,M}(x_k),\qquad
 q\longmapsto\gamma_{k,j}\star q.
\end{equation}
Definimos el estado del campo y su actualización por
\begin{align}
 \mathsf T_{k,j}
 &:=
 \left(\mathsf T(x_j),
   (\operatorname{pref}_{\gamma_{k,j},M})_{M\geq j}\right),
 \label{eq:pdfv2-terminal-estado-t}\\
 \operatorname{Upd}^{\mathsf T}_{\varepsilon}\mathsf T_{k,j}
 &:=
 \left(\mathsf T(x_{j+1}),
   (\operatorname{pref}_{\gamma_{k,j+1},M})_{M\geq j+1}\right),
 \qquad \gamma_{k,j+1}:=\gamma_{k,j}\star\varepsilon.
 \label{eq:pdfv2-terminal-update-t-explicita}
\end{align}
En particular, \(\mathsf T_{k,k}\) usa
\(\gamma_{k,k}=I_{x_k}\). Los testigos
\eqref{eq:pdfv2-terminal-testigo-neutro} y la sobreyectividad de \(\pi\)
prueban que el sistema es no vacío; nivel por nivel,
\begin{equation}
\label{eq:pdfv2-terminal-prefijo-natural}
 \pi^{\mathscr H}_{M+1,M}
 \operatorname{pref}_{\gamma_{k,j},M+1}
 =\operatorname{pref}_{\gamma_{k,j},M}
  \pi^{\mathscr H}_{M+1,M},
\end{equation}
porque ambos recorridos sólo eliminan la última emisión de \(q\).

\paragraph{13. Lector interno \(\mathsf L\).}
El lector no contiene ningún objeto clásico. Su tipo es la categoría de
caminos del único constructor dependiente
\(\mathfrak D_{j,r}\) de \eqref{eq:pdfv2-cor-constructor-unico}. Ponemos
\begin{equation}
\label{eq:pdfv2-terminal-transicion-lector-tipado}
 u^{\mathfrak D}_{j+1,j}
 :=u_{(j+1,r),(j,r)}:
 \mathfrak D_{j+1,r}(x_{j+1})
 \longrightarrow\mathfrak D_{j,r}(x_j),
\end{equation}
la componente del transporte común
\eqref{eq:pdfv2-cor-naturalidad-unica} que trunca la última emisión, reduce
los coeficientes y precompone los puertos. Por composición,
\begin{equation}
\label{eq:pdfv2-terminal-transicion-lector-composicion}
 u^{\mathfrak D}_{m,k}
 =u^{\mathfrak D}_{\ell,k}u^{\mathfrak D}_{m,\ell}
 \qquad(m\geq\ell\geq k).
\end{equation}
\begin{align}
 \mathsf L_{k,k}
 &:=(\mathfrak D_{0,r}(x_0))\star
 \mathop{\star}_{i=1}^{k}
 \bigl(\mathfrak D_{i,r}(x_i),u^{\mathfrak D}_{i,i-1},
       \mathcal O(\mathfrak e_{\varepsilon_i})\bigr),
 \label{eq:pdfv2-terminal-semilla-l}\\
 \operatorname{Upd}^{\mathsf L}_{\varepsilon}\mathsf L_{k,j}
 &:=
 \mathsf L_{k,j}\star
 \bigl(\mathfrak D_{j+1,r}(x_{j+1}),
 u^{\mathfrak D}_{j+1,j},\mathcal O(\mathfrak e_\varepsilon)\bigr).
 \label{eq:pdfv2-terminal-update-l}
\end{align}
La totalidad de \(\mathfrak D\) y
\eqref{eq:pdfv2-terminal-transicion-lector-composicion} prueban no vacuidad
y naturalidad.

Las trece semillas anteriores forman un único elemento
\begin{equation}
\label{eq:pdfv2-terminal-semilla-trece}
 \mathfrak r^0_{k,r}:=
 \langle
 \mathsf F_{k,k};\mathsf E_{k,k};\mathsf R_{k,k};\mathsf N_{k,k};
 \mathsf O_{k,k};\mathsf K_{k,k};\mathsf M_{k,k};\mathsf B_{k,k};
 \mathsf G_{k,k};\mathsf X_{k,k};\mathsf S_{k,k};\mathsf T_{k,k};
 \mathsf L_{k,k}
 \rangle\in\mathfrak R_{k,r}.
\end{equation}
La acción de una arista es la aplicación total
\begin{equation}
\label{eq:pdfv2-terminal-actualizacion-unica}
 \operatorname{Upd}_\varepsilon:
 \mathfrak R_{j,r}(\gamma_{0,j})
 \longrightarrow\mathfrak R_{j+1,r}(\gamma_{0,j+1}),\qquad
 \operatorname{Upd}_\varepsilon:=
 \langle\operatorname{Upd}^{\mathsf F}_\varepsilon,\ldots,
 \operatorname{Upd}^{\mathsf L}_\varepsilon\rangle,
\end{equation}
con las trece acciones
\eqref{eq:pdfv2-terminal-update-f}--
\eqref{eq:pdfv2-terminal-update-l}. Finalmente, el registro terminal se
\emph{genera} por el fold
\begin{equation}
\label{eq:pdfv2-terminal-registro-trece}
 \boxed{
 \mathfrak r_{k,N}(h):=
 \operatorname{Upd}_{\varepsilon_N}\circ\cdots\circ
 \operatorname{Upd}_{\varepsilon_{k+1}}(\mathfrak r^0_{k,r}).}
\end{equation}
No es el desempaquetado de una copia intacta del estado.
Sobre la imagen de este fold se define, componente a componente, el
truncamiento
\begin{equation}
\label{eq:pdfv2-terminal-truncamiento-trece-def}
 \operatorname{tr}_{13;j+1,j}
 \mathfrak r_{k,j+1}(h\star\varepsilon)
 :=\mathfrak r_{k,j}(h).
\end{equation}
Está bien definido: los campos \(\mathsf E\) y \(\mathsf G\) conservan la
lista ordenada de emisiones y prefijos, de modo que la igualdad de dos
registros implica la igualdad de sus historias etiquetadas. En
\(\mathsf F,\mathsf E,\mathsf R,\mathsf N,\mathsf G,\mathsf T,\mathsf L\)
se elimina la última entrada o bloque; en \(\mathsf O,\mathsf K,\mathsf M,
\mathsf B\) se aplica la fórmula acumulativa al prefijo recuperado; y en
\(\mathsf X,\mathsf S\) se restringe el cono de colas. No se invierte una
unión idempotente ni se selecciona una continuación.

\begin{proposition}[Totalidad y naturalidad de las trece actualizaciones]
\label{prop:pdfv2-terminal-trece-total-natural}
Todas las semillas de \eqref{eq:pdfv2-terminal-semilla-trece} existen, cada
\(\operatorname{Upd}_\varepsilon\) es total y, para toda simetría coherente
\(a\) y todo truncamiento de horizonte,
\begin{equation}
\label{eq:pdfv2-terminal-trece-naturalidad}
 a_*\operatorname{Upd}_\varepsilon
 =\operatorname{Upd}_{a\varepsilon}a_*,
 \qquad
 \operatorname{tr}_{13;j+1,j}
 \operatorname{Upd}_\varepsilon\mathfrak r_{k,j}(h)
 =\mathfrak r_{k,j}(h).
\end{equation}
\end{proposition}

\begin{proof}
La existencia de cada semilla fue exhibida en los trece párrafos. En
\(\mathsf F,\mathsf E,\mathsf G,\mathsf X,\mathsf S,\mathsf T,\mathsf L\),
la fórmula es concatenación con un mapa ya tipado. En
\(\mathsf R,\mathsf N,\mathsf O,\mathsf K,\mathsf M,\mathsf B\), la
fórmula es, respectivamente, concatenación de un bloque de ecuaciones
etiquetado, adjunción de un dato numérico, transporte de hojas, composición
de cociclos, acción afín y mapa de cadena. Las identidades de naturalidad
citadas después de cada fórmula prueban la primera igualdad de
\eqref{eq:pdfv2-terminal-trece-naturalidad} componente a componente. La
segunda es exactamente
\eqref{eq:pdfv2-terminal-truncamiento-trece-def}; su buena definición usa la
historia ordenada conservada por \(\mathsf E\) y \(\mathsf G\).
\end{proof}

\section{Terminal multivaluado sin selector}

El terminal finito generado por \(\widehat x_k\) es la multisecci\'on
\begin{equation}
\label{eq:pdfv2-terminal-multivaluado}
 \mathfrak T_{k,N}(\widehat x_k)
 :=\left\{\mathfrak r_{k,N}(h):
            h\in\mathscr H_{k,N}(\widehat x_k)\right\}.
\end{equation}
La expresi\'on de la derecha devuelve el conjunto entero. No contiene una
regla que distinga una historia privilegiada.

El ambiente terminal de nivel se obtiene reuniendo esas multisecciones sin
identificar registros:
\begin{equation}
\label{eq:pdfv2-terminal-ambiente-nivel}
 \mathfrak T_{k,N}
 :=\bigcup_{\widehat x_k\in\widehat{\mathcal X}^{\rm tot}_k}
   \mathfrak T_{k,N}(\widehat x_k).
\end{equation}

La versi\'on lineal, cuando se necesita sumar contribuciones, se define en
el m\'odulo libre por
\begin{equation}
\label{eq:pdfv2-terminal-evaluacion-todas}
 \begin{aligned}
 \mathsf E_{k,N}^{\rm all}:
 \mathbb Z[\widehat{\mathcal X}^{\rm tot}_k]&\longrightarrow
 \mathbb Z[\mathfrak T_{k,N}],\\
 [\widehat x_k]&\longmapsto
 \sum_{h\in\mathscr H_{k,N}(\widehat x_k)}
 [\mathfrak r_{k,N}(h)].
 \end{aligned}
\end{equation}
Cada testigo aparece con su multiplicidad. Si existe una medida proyectiva
positiva \(\mu\), la versi\'on isom\'etrica usa todas las probabilidades
condicionales:
\begin{equation}
\label{eq:pdfv2-terminal-evaluacion-ponderada}
 \widehat{\mathsf E}_{k,N}^{\mu}e_{\widehat x_k}
 :=\sum_{h\in\mathscr H_{k,N}(\widehat x_k)}
 \mu_{\widehat x_k,N}(h)^{1/2}
 e_{\mathfrak r_{k,N}(h)}.
\end{equation}
La proyectividad da
\begin{equation}
\label{eq:pdfv2-terminal-isometria}
 \left\|\widehat{\mathsf E}_{k,N}^{\mu}e_{\widehat x_k}\right\|^2
 =\sum_{h\in\mathscr H_{k,N}(\widehat x_k)}
   \mu_{\widehat x_k,N}(h)=1.
\end{equation}
La suma no se sustituye por ``la historia que genera el cilindro''.

\section{La conexi\'on non\'adica como correspondencia}

La holonom\'ia non\'adica fue generada por el od\'ometro y la conexi\'on TPK,
sin hip\'otesis de pertenencia, en
\eqref{eq:continuo-comun-gamma9-apice}--
\eqref{eq:continuo-comun-gamma9-relacion}. En este cap\'itulo usamos
\begin{equation}
\label{eq:pdfv2-terminal-apice-gamma9}
 \mathsf Z_{9,k}:=\mathsf Z^{\Gamma}_{9,k}.
\end{equation}
Los mapas
\begin{equation}
\label{eq:pdfv2-terminal-span-gamma9}
 \widehat{\mathcal X}^{\rm tot}_k
 \xleftarrow{\ s_{9,k}\ }
 \mathsf Z^{\Gamma}_{9,k}
 \xrightarrow{\ t_{9,k}\ }
 \widehat{\mathcal X}^{\rm tot}_{k+9}
\end{equation}
son los de \eqref{eq:continuo-comun-gamma9-span}; su fuente es sobreyectiva
por \cref{prop:continuo-comun-gamma9-total-natural}. Cada punto del \`apice
conserva las nueve emisiones fase--compatibles. No se reemplaza por el par
de sus extremos.

Para \(N\geq k+9\), definimos el \`apice de horizonte antes de linealizar:
\begin{equation}
\label{eq:pdfv2-terminal-apice-horizonte-gamma9}
 \mathsf Z^\Gamma_{9;k,N}:=
 \coprod_{z=(\widehat x_k,\boldsymbol\omega)
            \in\mathsf Z^\Gamma_{9,k}}
 \operatorname{Term}_{k+9,N}(t_{9,k}z).
\end{equation}
Un elemento se escribe \((z,q)\), donde \(\boldsymbol\omega\) es la palabra
non\'adica de nueve emisiones y \(q\) su cola hasta \(N\). Los mapas sobre
historias usan el codominio etiquetado
\begin{equation}
\label{eq:pdfv2-terminal-gamma9-codominio-colas}
 \mathscr D^\Gamma_{9;k,N}:=
 \coprod_{z\in\mathsf Z^\Gamma_{9,k}}
 \operatorname{Term}_{k+9,N}(t_{9,k}z).
\end{equation}
Así los mapas son
\begin{align}
 S_{k,N}(z,q)&:=\boldsymbol\omega\star q,
 \label{eq:pdfv2-terminal-gamma9-s-horizonte}\\
 T_{k,N}(z,q)&:=(z,q)\in\mathscr D^\Gamma_{9;k,N}.
 \label{eq:pdfv2-terminal-gamma9-t-horizonte}
\end{align}
El primero aterriza en
\(\coprod_x\mathscr H_{k,N}(x)\); el segundo conserva expresamente el
sumando etiquetado por \(z\).

El mapa de enlace del \`apice es la acci\'on expl\'icita
\begin{equation}
\label{eq:pdfv2-terminal-gamma9-enlace-apice}
 \pi^Z_{N+1,N}(z,q\star\varepsilon):=(z,q).
\end{equation}
Es sobreyectivo: si \((z,q)\) est\'a dado, se a\~nade a la cola la emisi\'on
neutra del extremo de \(q\). El codominio de colas tiene el enlace
\begin{equation}
\label{eq:pdfv2-terminal-gamma9-enlace-codominio}
 \pi^{\mathscr D}_{N+1,N}(z,q\star\varepsilon):=(z,q).
\end{equation}
Denotamos por
\(\pi^{\mathscr H,j}_{N+1,N}\) el mapa
\eqref{eq:pdfv2-terminal-truncamiento-historias} con nivel inicial \(j\).
Entonces
\begin{align}
 \pi^{\mathscr H,k}_{N+1,N}S_{k,N+1}
 &=S_{k,N}\pi^Z_{N+1,N},
 \label{eq:pdfv2-terminal-gamma9-cuadrado-fuente}\\
 \pi^{\mathscr D}_{N+1,N}T_{k,N+1}
 &=T_{k,N}\pi^Z_{N+1,N}.
 \label{eq:pdfv2-terminal-gamma9-cuadrado-destino}
\end{align}
En el primer cuadrado ambos recorridos env\'ian
\(\boldsymbol\omega\star q\star\varepsilon\) a
\(\boldsymbol\omega\star q\); en el segundo env\'ian
\((z,q\star\varepsilon)\) a \((z,q)\). Por descomposici\'on \`unica de una palabra
en sus primeras nueve emisiones y su cola, \(S_{k,N}\) es biyectivo. El
segundo cuadrado es cartesiano por la definici\'on coproducto--fibrada de
\eqref{eq:pdfv2-terminal-apice-horizonte-gamma9}. Son los cuadrados de
Beck--Chevalley que conservan testigos y multiplicidades.

El \`apice correspondiente de registros, ahora generado por las trece
actualizaciones, es
\begin{equation}
\label{eq:pdfv2-terminal-apice-terminal-gamma9}
 \begin{split}
 \mathsf Z^{\mathfrak T}_{9;k,N}:=\bigl\{
 (&z,q;
   \mathfrak r_{k,N}(\boldsymbol\omega\star q),
   \mathfrak r_{k+9,N}(q)):\\
 &z=(\widehat x_k,\boldsymbol\omega)\in
       \mathsf Z^\Gamma_{9,k},\quad
 q\in\operatorname{Term}_{k+9,N}(t_{9,k}z)
 \bigr\}.
 \end{split}
\end{equation}
Sus mapas fuente y destino leen, respectivamente, el primer o el segundo
registro, pero nunca eliminan de su dominio el testigo \((z,q)\):
\begin{equation}
\label{eq:pdfv2-terminal-span-terminal}
 \mathfrak T_{k,N}
 \xleftarrow{\ s^{\mathfrak T}_{9;k,N}\ }
 \mathsf Z^{\mathfrak T}_{9;k,N}
 \xrightarrow{\ t^{\mathfrak T}_{9;k,N}\ }
 \mathfrak T_{k+9,N}.
\end{equation}
Si \(\zeta=(z,q;r_s,r_t)\), su acción es
\begin{equation}
\label{eq:pdfv2-terminal-span-terminal-accion}
 s^{\mathfrak T}_{9;k,N}(\zeta):=r_s,
 \qquad
 t^{\mathfrak T}_{9;k,N}(\zeta):=r_t.
\end{equation}
Su enlace se define por
\begin{equation}
\label{eq:pdfv2-terminal-gamma9-enlace-registros}
 \begin{split}
 \pi^{Z,\mathfrak T}_{N+1,N}
 \bigl(&z,q\star\varepsilon;
       \mathfrak r_{k,N+1}
          (\boldsymbol\omega\star q\star\varepsilon),
       \mathfrak r_{k+9,N+1}(q\star\varepsilon)\bigr)\\
 &:=(z,q;\mathfrak r_{k,N}(\boldsymbol\omega\star q),
       \mathfrak r_{k+9,N}(q)).
 \end{split}
\end{equation}
La pareja \((z,q)\) es una coordenada del \`apice, no una representación
elegida de una terna; por ello el enlace está bien definido incluso si dos
registros terminales coinciden. La compatibilidad de
\(\operatorname{Upd}\) con el borrado de la \`ultima arista prueba que
\eqref{eq:pdfv2-terminal-gamma9-enlace-registros} conmuta con ambos mapas de
\eqref{eq:pdfv2-terminal-span-terminal}.

Si se desea una acci\'on sobre m\'odulos libres, se linealiza esta
correspondencia sumando todos los testigos:
\begin{equation}
\label{eq:pdfv2-terminal-linealizacion-span}
 [\Gamma_9]_{k,N}[r]
 :=\sum_{\zeta\in(s^{\mathfrak T}_{9;k,N})^{-1}(r)}
 [t^{\mathfrak T}_{9;k,N}(\zeta)].
\end{equation}
La f\'ormula conserva multiplicidades. La igualdad de fases despu\'es de
nueve pasos no contrae los estados, los testigos ni sus ledgers.

\section{Una sola ley de naturalidad}

Para \(k\leq\ell\leq N\), toda historia de
\(\mathscr H_{k,N}(\widehat x_k)\) posee un \'unico prefijo
\(p\in\mathscr P_{\ell\mid k}(\widehat x_k)\). Por tanto hay una descomposici\'on
disjunta por prefijos, expresada sin elegir ninguno:
\begin{equation}
\label{eq:pdfv2-terminal-particion-historias}
 \mathscr H_{k,N}(\widehat x_k)
 =\bigsqcup_{p\in\mathscr P_{\ell\mid k}(\widehat x_k)}
 \left\{p\star q:q\in\mathscr H_{\ell,N}(p)\right\}.
\end{equation}
Sea \(\operatorname{Conc}_p\) la operaci\'on que antepone el registro del
prefijo \(p\) y aplica sucesivamente
\eqref{eq:pdfv2-terminal-actualizacion-unica}. Definimos la transici\'on
com\'un sobre una familia completa de terminales por
\begin{equation}
\label{eq:pdfv2-terminal-transicion-comun}
 \mathfrak u_{\ell,k}
 \left(
   \bigl(\mathfrak T_{\ell,N}(p)\bigr)_{p\in\mathscr P_{\ell\mid k}(\widehat x_k)}
 \right)
 :=\bigcup_{p\in\mathscr P_{\ell\mid k}(\widehat x_k)}
   \operatorname{Conc}_p\mathfrak T_{\ell,N}(p).
\end{equation}

Escribiendo
\(\mathbf G_{k,N}(\widehat x_k):=\mathfrak T_{k,N}(\widehat x_k)\), toda la naturalidad se
concentra en la identidad \'unica
\begin{equation}
\label{eq:pdfv2-terminal-naturalidad-unica}
 \boxed{
 \mathfrak u_{\ell,k}
 \left(
   \bigl(\mathbf G_{\ell,N}(p)\bigr)_{p\in\mathscr P_{\ell\mid k}(\widehat x_k)}
 \right)
 =\mathbf G_{k,N}(\widehat x_k).}
\end{equation}
No hay una ley distinta para cada discriminante interno. Las cinco
caracter\'isticas obedecen
\eqref{eq:pdfv2-terminal-naturalidad-unica} porque cada una se actualiza
mediante el mismo prefijo, la misma extensi\'on y el mismo ledger.

La correspondencia non\'adica es compatible con esta identidad mediante los
mapas ya construidos, no mediante una condici\'on impl\'icita. Para todo
horizonte, \eqref{eq:pdfv2-terminal-gamma9-cuadrado-fuente} y
\eqref{eq:pdfv2-terminal-gamma9-cuadrado-destino} conservan el testigo
non\'adico con su multiplicidad, y
\eqref{eq:pdfv2-terminal-gamma9-enlace-registros} hace lo mismo en las
trece coordenadas. Si \(a\) es una simetr\'ia coherente del propietario
com\'un, su acci\'on
\[
 a^Z(z,q):=(a^Zz,aq)
\]
satisface
\begin{equation}
\label{eq:pdfv2-terminal-gamma9-equivarianza-completa}
 S a^Z=aS,\qquad T a^Z=aT,\qquad
 \pi^Za^Z=a^Z\pi^Z,\qquad
 s^{\mathfrak T}a^Z=as^{\mathfrak T},\qquad
 t^{\mathfrak T}a^Z=at^{\mathfrak T}.
\end{equation}
Estas igualdades se verifican entrada por entrada usando
\eqref{eq:continuo-comun-gamma9-naturalidad} y la naturalidad de las trece
actualizaciones. Son la condici\'on material que permite usar
\eqref{eq:pdfv2-terminal-linealizacion-span}; una mera igualdad de valores
visibles no la sustituye.

\section{Paso al l\'imite sin intercambiar imagen y l\'imite}

Truncar una historia y cada uno de los trece campos define
\begin{equation}
\label{eq:pdfv2-terminal-transicion-horizonte}
 \pi^{\mathfrak T}_{N+1,N}:
 \mathfrak T_{k,N+1}(\widehat x_k)\longrightarrow
 \mathfrak T_{k,N}(\widehat x_k),
 \qquad
 \pi^{\mathfrak T}_{N+1,N}\mathfrak r_{k,N+1}(h)
 =\mathfrak r_{k,N}(h_{\leq N}).
\end{equation}
Es sobreyectiva por
\eqref{eq:pdfv2-terminal-truncamiento-historias}. El terminal completo es
entonces
\begin{equation}
\label{eq:pdfv2-terminal-limite-inverso}
 \mathfrak T_{k,\infty}(\widehat x_k)
 :=\left\{(r_N)_{N\geq k}:
 r_N\in\mathfrak T_{k,N}(\widehat x_k),\quad
 \pi^{\mathfrak T}_{N+1,N}r_{N+1}=r_N\right\}.
\end{equation}
Los conjuntos finitos de
\eqref{eq:pdfv2-terminal-multivaluado} son no vac\'ios y sus transiciones
son sobreyectivas; por ello
\begin{equation}
\label{eq:pdfv2-terminal-limite-no-vacio}
 \mathfrak T_{k,\infty}(\widehat x_k)\neq\varnothing.
\end{equation}

El \`apice non\'adico en el l\'imite se define, no se infiere de una imagen:
\begin{equation}
\label{eq:pdfv2-terminal-gamma9-apice-limite-raw}
 \mathsf Z^\Gamma_{9;k,\infty}
 :=\varprojlim_{N\geq k+9}
 \bigl(\mathsf Z^\Gamma_{9;k,N},\pi^Z_{N+1,N}\bigr).
\end{equation}
Como los enlaces no modifican el testigo finito \(z\) y s\'olo truncan su
cola, existe el isomorfismo expl\'icito
\begin{equation}
\label{eq:pdfv2-terminal-gamma9-apice-limite-msect}
 \mathsf Z^\Gamma_{9;k,\infty}
 \xrightarrow{\ \cong\ }
 \coprod_{z\in\mathsf Z^\Gamma_{9,k}}
 \operatorname{Msect}(t_{9,k}z),\qquad
 ((z,q_N))_N\longmapsto(z,(q_N)_N).
\end{equation}
La inversa toma un testigo \(z\) y una cola compatible \((q_N)_N\), y forma
\(((z,q_N))_N\). Ambas composiciones son identidades coordenada a
coordenada. Los mapas l\'imite son
\begin{equation}
\label{eq:pdfv2-terminal-gamma9-mapas-limite-raw}
 S_{k,\infty}((z_N)_N):=(S_{k,N}z_N)_N,\qquad
 T_{k,\infty}((z_N)_N):=(T_{k,N}z_N)_N;
\end{equation}
est\'an bien definidos precisamente por
\eqref{eq:pdfv2-terminal-gamma9-cuadrado-fuente} y
\eqref{eq:pdfv2-terminal-gamma9-cuadrado-destino}.

Del mismo modo se define el \`apice de registros
\begin{equation}
\label{eq:pdfv2-terminal-gamma9-apice-registros-limite}
 \mathsf Z^{\mathfrak T}_{9;k,\infty}
 :=\varprojlim_{N\geq k+9}
 \bigl(\mathsf Z^{\mathfrak T}_{9;k,N},
       \pi^{Z,\mathfrak T}_{N+1,N}\bigr).
\end{equation}
Las correspondencias
\eqref{eq:pdfv2-terminal-span-terminal} forman entonces, componente a
componente, el span
\begin{equation}
\label{eq:pdfv2-terminal-span-limite}
 \mathfrak T_{k,\infty}
 \xleftarrow{\ s^{\mathfrak T}_{9;k,\infty}\ }
 \mathsf Z^{\mathfrak T}_{9;k,\infty}
 \xrightarrow{\ t^{\mathfrak T}_{9;k,\infty}\ }
 \mathfrak T_{k+9,\infty}.
\end{equation}
No se ha identificado la imagen de un l\'imite con el l\'imite de unas
im\'agenes: \eqref{eq:pdfv2-terminal-limite-inverso},
\eqref{eq:pdfv2-terminal-gamma9-apice-limite-raw} y
\eqref{eq:pdfv2-terminal-gamma9-apice-registros-limite} construyen
directamente las tres familias compatibles y sus mapas.

\section{Recuperaci\'on interna de las trece coordenadas}

La recuperaci\'on no desempaqueta nombres: repite el fold que gener\'o el
registro. Para
\(h=(\varepsilon_{k+1},\ldots,\varepsilon_N)\), definimos
\begin{equation}
\label{eq:pdfv2-terminal-recuperacion-trece}
 \boxed{
 \operatorname{Rec}_{13;k,N}(h):=
 \operatorname{Upd}_{\varepsilon_N}\circ\cdots\circ
 \operatorname{Upd}_{\varepsilon_{k+1}}
 (\mathfrak r^0_{k,r}).}
\end{equation}
Por \eqref{eq:pdfv2-terminal-registro-trece},
\(\operatorname{Rec}_{13;k,N}(h)=\mathfrak r_{k,N}(h)\), pero la igualdad
es ahora una identidad entre dos composiciones de acciones, no la
proyecci\'on de una entrada preservada. Si
\(\operatorname{pr}_{\mathsf A}\) denota el proyector del registro al campo
\(\mathsf A\in\{\mathsf F,\mathsf E,\mathsf R,\mathsf N,\mathsf O,
\mathsf K,\mathsf M,\mathsf B,\mathsf G,\mathsf X,\mathsf S,\mathsf T,
\mathsf L\}\), entonces
\begin{equation}
\label{eq:pdfv2-terminal-recuperacion-campo-fold}
 \operatorname{pr}_{\mathsf A}\operatorname{Rec}_{13;k,N}(h)
 =\operatorname{Upd}^{\mathsf A}_{\varepsilon_N}\circ\cdots\circ
  \operatorname{Upd}^{\mathsf A}_{\varepsilon_{k+1}}
  (\mathsf A_{k,k}).
\end{equation}
La prueba es inducci\'on en \(N-k\): para \(N=k\) ambos lados son la
semilla; el paso inductivo aplica la componente \(\mathsf A\) de
\(\operatorname{Upd}_{\varepsilon_N}\).

Adem\'as,
\begin{equation}
\label{eq:pdfv2-terminal-recuperacion-compatible}
 \operatorname{Rec}_{13;k,N}(h_{\leq N})
 =\operatorname{tr}_{13,N+1,N}
   \bigl(\operatorname{Rec}_{13;k,N+1}(h)\bigr),
\end{equation}
donde \(\operatorname{tr}_{13,N+1,N}\) elimina s\'olo la \'ultima
actualizaci\'on y restringe su cono de colas. En consecuencia, un elemento
de \(\mathfrak T_{k,\infty}(\widehat x_k)\) determina una familia compatible de las
trece coordenadas a toda profundidad.

\section{Teorema terminal conjunto}

\begin{theorem}[Terminal multivaluado, naturalidad y l\'imite del continuo]
\label{thm:pdfv2-terminal-conjunto}
Para todo \(\widehat x_k\in\widehat{\mathcal X}^{\rm tot}_k\), las reglas
\eqref{eq:pdfv2-terminal-historias-finita}--
\eqref{eq:pdfv2-terminal-recuperacion-compatible} construyen, a partir del
\'unico continuo discreto:
\begin{enumerate}
 \item una multisecci\'on terminal finita no vac\'ia en cada horizonte;
 \item una evaluaci\'on que conserva todas las historias y multiplicidades;
 \item la conexi\'on non\'adica como correspondencia con testigos;
 \item la identidad \'unica de naturalidad
       \eqref{eq:pdfv2-terminal-naturalidad-unica};
 \item un terminal inverso no vac\'io y una correspondencia non\'adica en el
       l\'imite;
 \item la recuperaci\'on compatible de fibra, extensiones, relaciones,
       n\'umeros, orientaci\'on, carry, memoria, frontera, ruta,
       multisecci\'on, supervivencia, terminal y lector.
\end{enumerate}
Las cinco caracter\'isticas intr\'insecas permanecen correlacionadas dentro
de cada registro y de cada testigo de extensi\'on.
\end{theorem}

\begin{proof}
La no vacuidad de
\(\mathscr H_{k,N}(\widehat x_k)\) es no vac\'ia por
\cref{prop:continuo-comun-arbol-finito-no-vacio}, y su truncamiento es
sobreyectivo por \cref{lem:continuo-comun-terminal-sobreyectivo}. La regla
\eqref{eq:pdfv2-terminal-actualizacion-unica} aplica las identidades APP,
TRIT y TPK sobre cada testigo \(\varepsilon_j\); por inducci\'on en
\(N-k\), produce todos los campos de
\eqref{eq:pdfv2-terminal-registro-trece} y conserva sus relaciones.

Tomar el conjunto completo en
\eqref{eq:pdfv2-terminal-multivaluado}, o la suma completa en
\eqref{eq:pdfv2-terminal-evaluacion-todas}, demuestra que el terminal no
contiene un selector. La versi\'on ponderada es isom\'etrica por la igualdad
proyectiva \eqref{eq:pdfv2-terminal-isometria}.

El testigo \(z=(\widehat x_k,\boldsymbol\omega)\) lleva la acción de fase de cada una
de sus nueve emisiones; por
\cref{prop:continuo-comun-gamma9-retorno-memoria} define los dos mapas de
\eqref{eq:pdfv2-terminal-span-gamma9} sin convertir la relación en una
función ni equiparar fase y estado. Las fórmulas
\eqref{eq:pdfv2-terminal-gamma9-s-horizonte}--
\eqref{eq:pdfv2-terminal-gamma9-enlace-registros} prueban los dos cuadrados
de horizonte y la misma construcción sobre registros. Sumar el espacio
fuente de cada testigo da
\eqref{eq:pdfv2-terminal-linealizacion-span} sin borrar multiplicidades.

La partici\'on
\eqref{eq:pdfv2-terminal-particion-historias} es disjunta porque toda
historia tiene un prefijo de longitud \(\ell\) determinado. Concatenar cada
prefijo con todas sus continuaciones reconstruye una vez y s\'olo una cada
historia de \(\mathscr H_{k,N}(\widehat x_k)\). Esto prueba directamente
\eqref{eq:pdfv2-terminal-naturalidad-unica} para el registro entero; no se
necesitan cinco demostraciones inconexas.

La sobreyectividad de
\eqref{eq:pdfv2-terminal-transicion-horizonte} y la finitud de cada nivel
dan la no vacuidad de
\eqref{eq:pdfv2-terminal-limite-inverso}. Los enlaces explícitos
\eqref{eq:pdfv2-terminal-gamma9-enlace-apice} y
\eqref{eq:pdfv2-terminal-gamma9-enlace-registros}, junto con sus cuadrados,
construyen los límites
\eqref{eq:pdfv2-terminal-gamma9-apice-limite-raw} y
\eqref{eq:pdfv2-terminal-gamma9-apice-registros-limite}; de ellos sale
\eqref{eq:pdfv2-terminal-span-limite}. Finalmente,
\eqref{eq:pdfv2-terminal-recuperacion-trece} vuelve a ejecutar las
actualizaciones sucesivas, y
\eqref{eq:pdfv2-terminal-recuperacion-compatible} demuestra que esas
lecturas pasan conjuntamente al l\'imite.
\end{proof}

\section{Falsadores materiales}

\begin{proposition}[Falsadores del terminal conjunto]
\label{prop:pdfv2-terminal-falsadores}
Cualquiera de las operaciones siguientes sustituye el objeto construido y
bloquea la conclusi\'on del
\cref{thm:pdfv2-terminal-conjunto}:
\begin{enumerate}
 \item enviar un cilindro a una sola historia cuando admite varias;
 \item elegir una cola del campo terminal \(\mathsf T\) y eliminar las
       restantes;
 \item reemplazar \(\Gamma_9\) por una funci\'on sin demostrar unicidad de
       cada destino y de su testigo;
 \item linealizar una correspondencia ignorando multiplicidades;
 \item identificar retorno de fase con retorno de estado;
 \item conservar la entrada intacta como primera coordenada y usar su
       recuperabilidad como supuesta generaci\'on;
 \item omitir cualquiera de los trece campos de
       \eqref{eq:pdfv2-terminal-registro-trece};
 \item permitir que dos caracter\'isticas intr\'insecas usen historias,
       carries, fronteras o terminales diferentes;
 \item sustituir
       \eqref{eq:pdfv2-terminal-naturalidad-unica} por cinco afirmaciones
       independientes;
 \item afirmar compatibilidad operatorial de \(\Gamma_9\) sin conservar el
       espacio de testigos y sus multiplicidades;
 \item intercambiar, sin prueba, formaci\'on de im\'agenes y paso al l\'imite;
 \item presentar una lista de trece nombres sin las acciones
       \eqref{eq:pdfv2-terminal-actualizacion-unica} que los producen.
\end{enumerate}
\end{proposition}

\begin{proof}
Los cinco primeros actos contraen la multisecci\'on o la correspondencia. Los
cuatro siguientes rompen el soporte y la naturalidad comunes. Los tres
\'ultimos eliminan la prueba generativa o alteran el paso al l\'imite. En
cada caso, el objeto resultante deja de satisfacer al menos una de
\eqref{eq:pdfv2-terminal-multivaluado},
\eqref{eq:pdfv2-terminal-span-terminal},
\eqref{eq:pdfv2-terminal-naturalidad-unica} o
\eqref{eq:pdfv2-terminal-recuperacion-compatible}; por tanto, la conclusi\'on
no se transfiere.
\end{proof}

\begin{statusbox}
\texttt{FORMALIZACION\_NUEVA}: organizaci\'on del terminal completo, la
naturalidad \'unica y el paso al l\'imite.\quad
\texttt{RESULTADO\_RECUPERADO}: multisecci\'on superviviente, carries,
ledger, frontera, ruta y conexi\'on non\'adica.\quad
La linealización sobre módulos libres está demostrada por
\eqref{eq:pdfv2-terminal-linealizacion-span} y sus cuadrados de
Beck--Chevalley. Una promoción posterior a un operador acotado sobre una
completación analítica exige además declarar esa norma y probar la cota; la
correspondencia \eqref{eq:pdfv2-terminal-span-limite} no depende de esa
interfaz.
\end{statusbox}
